\par\Needspace{6\baselineskip}\medskip
\begin{center}\hypertarget{native-182-article136}{}\pdfbookmark[2]{136.}{article-136}136.\end{center}
Ante omnia tres locos heliocentricos terrae in sphaera coelesti $A,A'\;A''$
(fig. 4) cum tribus locis geocentricis respondentibus corporis coelestis $B,B',B''$
per circulos maximos iungere, atque tum positionem horum circulorum maximorum
respectu eclipticae (siquidem eclipticam pro plano fundamentali adoptamus),
tum situm punctorum $B,B',B''$ in ipsis computare oportet. Sint $\alpha,\alpha',\alpha''$
tres corporis coelestis longitudines geocentricae; $\beta,\beta',\beta''$ latitudines; $l,l',l''$
longitudines heliocentricae terrae, cuius latitudines statuimus $=0$ (art. 117, 72).
Sint porro $\gamma,\gamma',\gamma''$ circulorum maximorum ab $A,A',A''$ resp. ad $B,B',B''$ ductorum
inclinationes ad eclipticam: quas inclinationes, ut in ipsarum determinatione
normam fixam sequamur, perpetuo respectu eius eclipticae partis mensurabimus,
quae a punctis $A,A',A''$ secundum ordinem signorum sita est, ita
ut ipsarum magnitudo a $0$ usque ad $360^0$ numeretur, sive quod eodem redit,
in parte boreali a $0$ usque ad $180^0$, in australi a $0$ usque ad $-180^0$. Arcus
$AB,A'B',A''B''$, quos semper intra $0$ et $180^0$ statuere licebit, designamus per
$\delta,\delta',\delta''$. Ita pro determinatione ipsarum $\gamma,\delta$ habemus formulas
\[
\begin{aligned}
1)\quad&\operatorname{tang}\gamma=\frac{\operatorname{tang}\beta}{\sin(\alpha-l)}\\[.7ex]
2)\quad&\operatorname{tang}\delta=\frac{\operatorname{tang}(\alpha-l)}{\cos\gamma}\mathpunct{\raisebox{.45ex}{,}}
\end{aligned}
\]
quibus si placet ad calculi confirmationem adiici possunt sequentes:
\[
\sin\delta=\frac{\sin\beta}{\sin\gamma}\mathpunct{\raisebox{.45ex}{,}}\qquad
\cos\delta=\cos\beta\cos(\alpha-l).
\]
Pro determinandis $\gamma',\delta',\gamma'',\delta''$ manifesto formulae prorsus analogae habentur.
Quodsi simul fuerit $\beta=0$, $\alpha-l=0$ vel $=180^0$, i.\,e. si corpus coeleste simul
in oppositione vel coniunctione atque in ecliptica fuerit, $\gamma$ fieret indeterminata:
at supponemus, hunc casum in nulla trium observationum locum habere.

Si loco eclipticae aequator tamquam planum fundamentale adoptatum est,
ad positionem trium circulorum maximorum respectu aequatoris determinandam
praeter inclinationes insuper requirentur rectascensiones intersectionum cum
\hypertarget{native-183}{}aequatore: nec non praeter distantias punctorum $B,B',B''$ ab his intersectionibus
etiam distantias punctorum $A,A',A''$ ab iisdem computare oportebit.
Quae quum a problemate in art. 110 tractato pendeant, formularum evolutioni
hic non immoramur.
